
\documentclass[11pt,a4paper]{article}

% ─── Packages ───
\usepackage[utf8]{inputenc}
\usepackage[T1]{fontenc}
\usepackage{lmodern}
\usepackage[margin=1in]{geometry}
\usepackage{amsmath,amssymb,amsfonts}
\usepackage{hyperref}
\usepackage{xcolor}
\usepackage{listings}
\usepackage{booktabs}
\usepackage{enumitem}
\usepackage{float}
\usepackage{titlesec}

\hypersetup{
    colorlinks=true,
    linkcolor=blue!60!black,
    citecolor=green!40!black,
    urlcolor=cyan!50!black,
}

\lstset{
    language=Python,
    basicstyle=\ttfamily\footnotesize,
    keywordstyle=\color{blue},
    commentstyle=\color{gray},
    stringstyle=\color{red!70!black},
    breaklines=true,
    frame=single,
    numbers=left,
    numberstyle=\tiny\color{gray},
    showstringspaces=false,
    tabsize=4,
}

\definecolor{enhance}{RGB}{0,100,0}
\newcommand{\enhance}[1]{\textcolor{enhance}{\textbf{#1}}}
\newcommand{\file}[1]{\texttt{#1}}
\newcommand{\fn}[1]{\texttt{#1}}

\title{
    \textbf{DynamicScaler: Enhancement Analysis Report} \\[6pt]
    \large Opportunities for Improving Code, Logic, and Mathematics \\
    for Higher-Quality Panoramic Video Generation
}
\author{Auto-generated Codebase Analysis}
\date{\today}

\begin{document}
\maketitle

\begin{abstract}
This document provides a systematic analysis of the \textbf{DynamicScaler} codebase---a CVPR 2025 project for seamless and scalable panoramic video generation.
We identify concrete opportunities where code quality, algorithmic logic, and mathematical formulations can be enhanced to yield higher-quality results.
Each section describes the current implementation, explains the limitation, and proposes an enhancement with mathematical justification.
\end{abstract}

\tableofcontents
\newpage

%======================================================================
\section{Overview of the Codebase}
%======================================================================

DynamicScaler generates 360\textdegree{} panoramic videos by:
\begin{enumerate}
    \item Representing the panorama in equirectangular format.
    \item Extracting perspective views via gnomonic projection.
    \item Denoising each view using a pre-trained video diffusion model (VideoCrafter).
    \item Merging denoised views back into the equirectangular domain using shift-window strategies.
    \item Optionally upscaling via a second-pass denoise stage.
\end{enumerate}

The pipeline orchestrates three key sampling methods: sphere panorama shift-window denoising, normal plane shift-window denoising, and upscaled shift-window denoising.

%======================================================================
\section{Mathematics: Equirectangular Projection}
%======================================================================

\subsection{Current Implementation}

In \file{panorama\_tensor\_utils.py} and \file{ring\_panorama\_tensor\_utils.py}, the perspective--equirectangular mapping is implemented in \fn{\_get\_uv()}.
The coordinate grid is constructed as:

\begin{lstlisting}
x = torch.linspace(-width/2, width/2 - 1, steps=width)
y = torch.linspace(-height/2, height/2 - 1, steps=height)
\end{lstlisting}

\subsection{Issue: Asymmetric Pixel Grid}

The grid ranges from $[-W/2, \; W/2-1]$ instead of the half-pixel-centered symmetric grid $[-(W-1)/2, \; (W-1)/2]$.
For an even-width image (e.g.\ $W=512$), the grid spans $[-256, 255]$, introducing a \textbf{0.5-pixel bias} in the horizontal and vertical directions.
This asymmetry propagates through the rotation matrix and longitude/latitude conversion, causing a subtle sub-pixel misalignment in every view extraction and re-projection.

\subsection{Proposed Enhancement}

Replace with a properly centered grid:
\begin{equation}
    x_i = \frac{(2i - (W-1))}{2}, \quad i = 0, 1, \ldots, W-1
\end{equation}
\begin{lstlisting}
x = torch.linspace(-(width-1)/2, (width-1)/2, steps=width)
y = torch.linspace(-(height-1)/2, (height-1)/2, steps=height)
\end{lstlisting}

This ensures the grid is \textbf{symmetric about the optical center}, improving geometric accuracy of the gnomonic projection.

%======================================================================
\section{Mathematics: Longitude--Latitude Mapping Normalization}
%======================================================================

\subsection{Current Implementation}

After computing spherical coordinates $(\text{lon}, \text{lat})$, the mapping to equirectangular pixel coordinates $(u, v)$ is:
\begin{equation}
    u = \frac{\text{lon}}{2\pi} \cdot (W-1), \qquad
    v = \frac{\text{lat} + \pi/2}{\pi} \cdot (H-1)
\end{equation}

\subsection{Issue: Off-by-One in Wrap-Around}

Using $(W-1)$ instead of $W$ means that the leftmost column ($u=0$) maps to $\text{lon}=0$ and the rightmost column ($u=W-1$) maps to $\text{lon}=2\pi$ -- but $\text{lon}=0$ and $\text{lon}=2\pi$ are \emph{the same physical direction}.
This causes the rightmost column to be a duplicate of the leftmost, wasting one column's worth of resolution and creating potential seam artifacts.

\subsection{Proposed Enhancement}

For equirectangular panoramas with horizontal periodicity:
\begin{equation}
    u = \frac{\text{lon}}{2\pi} \cdot W
\end{equation}
This spreads $W$ pixels across $[0, 2\pi)$ without duplication.
The vertical mapping should remain with $(H-1)$ since the poles are genuine boundary conditions:
\begin{equation}
    v = \frac{\text{lat} + \pi/2}{\pi} \cdot (H-1)
\end{equation}

%======================================================================
\section{Mathematics: Nearest-Neighbor Sampling}
%======================================================================

\subsection{Current Implementation}

In \fn{\_sample\_equirect\_tensor\_nearest()}, the nearest-neighbor sampling uses \fn{torch.floor()}:
\begin{lstlisting}
u0 = torch.floor(u).long()
v0 = torch.floor(v).long()
\end{lstlisting}

\subsection{Issue: Biased Rounding}

\fn{floor()} always rounds \emph{downward}, introducing a systematic directional bias.
For truly nearest-neighbor sampling, \fn{torch.round()} should be used so that each pixel maps to its geometrically closest neighbor:

\begin{equation}
    u_{\text{nn}} = \operatorname{round}(u), \quad v_{\text{nn}} = \operatorname{round}(v)
\end{equation}

\subsection{Proposed Enhancement}
\begin{lstlisting}
u0 = torch.round(u).long()
v0 = torch.round(v).long()
\end{lstlisting}

This reduces average reprojection error from $\sim 0.5$ px to $\sim 0.25$ px (under uniform distribution of sub-pixel offsets), yielding sharper view reconstructions.

%======================================================================
\section{Mathematics: Bilinear Splatting in \fn{set\_view\_tensor\_bilinear()}}
%======================================================================

\subsection{Current Implementation}

The bilinear write-back in \file{panorama\_tensor\_utils.py} uses a nested Python loop over batch and channel dimensions:
\begin{lstlisting}
for b in range(B):
    for c in range(self.C):
        accumulator[c].index_add_(0, idx00, view_flat[b,c] * w00)
        ...
\end{lstlisting}

\subsection{Issue 1: Computational Inefficiency}

The $O(B \times C)$ Python loop is extremely slow.
With $B=1$ and $C=4$ (latent channels), this involves 16 \fn{index\_add\_} calls per pixel set.

\subsection{Issue 2: Missing Normalization for Overlapping Writes}

When multiple source pixels from the view map to the same equirectangular pixel, the accumulated value is divided by the sum of weights.
However, this \textbf{overwrites the existing panorama content entirely} in the accumulated region, rather than blending with existing values.

\subsection{Proposed Enhancement}

\begin{enumerate}
    \item \textbf{Vectorize} the scatter operation using batched \fn{scatter\_add\_()}:
\begin{lstlisting}
# Flatten across B,C for a single scatter_add_ call
view_flat_all = view_flat.reshape(-1, H*W)
# Weighted values for each corner
for idx, w in [(idx00,w00),(idx01,w01),(idx10,w10),(idx11,w11)]:
    accum.scatter_add_(1, idx.unsqueeze(0).expand(B*C,-1),
                       view_flat_all * w.unsqueeze(0))
\end{lstlisting}

    \item \textbf{Blend with existing content} by introducing a user-controlled blending factor $\alpha \in [0,1]$:
\begin{equation}
    P_{\text{new}}(u,v) = \alpha \cdot \frac{\sum_k w_k \cdot V_k}{\sum_k w_k} + (1-\alpha) \cdot P_{\text{old}}(u,v)
\end{equation}
\end{enumerate}

%======================================================================
\section{Mathematics: Re-noising Formulation}
%======================================================================

\subsection{Current Implementation}

In \file{scheduler.py}, the \fn{re\_noise()} method converts a sample $x_a$ at noise level $a$ to noise level $b$:

\begin{equation}
    x_b = \sqrt{\frac{\bar{\alpha}_b}{\bar{\alpha}_a}} \, x_a + \sqrt{1 - \frac{\bar{\alpha}_b}{\bar{\alpha}_a}} \, \tilde{\epsilon}
\end{equation}

\subsection{Issue: Information-Theoretic Sub-optimality}

This formula is the ``direct'' re-noising approach, which adds \emph{independent} Gaussian noise $\tilde{\epsilon}$.
It does not preserve as much structural information as the DDIM-inversion-based re-noising, because the newly added noise is uncorrelated with the original noise trajectory.

\subsection{Proposed Enhancement: DDIM-Inversion-based Re-noising}

Instead of independent noise injection, use DDIM inversion to estimate what the noise prediction would be at step $a$, then forward the DDIM process to step $b$:

\begin{align}
    \hat{\epsilon}_a &= \frac{x_a - \sqrt{\bar{\alpha}_a} \cdot \hat{x}_0}{\sqrt{1 - \bar{\alpha}_a}} \\
    x_b &= \sqrt{\bar{\alpha}_b} \cdot \hat{x}_0 + \sqrt{1 - \bar{\alpha}_b - \sigma_b^2} \cdot \hat{\epsilon}_a + \sigma_b \cdot \epsilon
\end{align}

where $\hat{x}_0$ can be estimated by running a single denoising step.
This preserves the noise trajectory and leads to more coherent transitions between noise levels, especially important in the multi-scale upscaling pipeline.

%======================================================================
\section{Logic: Overlap Ratio Scheduling}
%======================================================================

\subsection{Current Implementation}

In \file{gen\_pano\_360.py}, the temporal overlap ratio schedule is built via:
\begin{lstlisting}
overlap_ratio_list_1_f_org = [0.75, 0.5]
overlap_ratio_list_1_f = [
    overlap_ratio_list_1_f_org[
        i * len(overlap_ratio_list_1_f_org) // num_inference_steps
    ]
    for i in range(num_inference_steps)
]
\end{lstlisting}

\subsection{Issue: Abrupt Step-Function Transition}

With only two values $[0.75, 0.5]$, the overlap ratio drops discontinuously from 0.75 to 0.5 at the midpoint of the denoising schedule.
This causes an abrupt change in the effective receptive field overlap, which can lead to visible temporal seam artifacts at the transition boundary.

\subsection{Proposed Enhancement: Smooth Cosine Annealing}

Replace the step function with a smooth cosine annealing schedule:
\begin{equation}
    r(t) = r_{\max} - \frac{r_{\max} - r_{\min}}{2} \left(1 - \cos\left(\frac{\pi t}{T}\right)\right)
\end{equation}
where $r_{\max}=0.75$, $r_{\min}=0.5$, $t$ is the timestep index, and $T$ is the total number of inference steps.

\begin{lstlisting}
import math
overlap_ratio_list_f = [
    r_max - (r_max - r_min)/2 * (1 - math.cos(math.pi * i / T))
    for i in range(T)
]
\end{lstlisting}

This provides a gradual transition, reducing visible boundary effects.

%======================================================================
\section{Logic: Denoised Overlap Merging Ratio Decay}
%======================================================================

\subsection{Current Implementation}

The \fn{merge\_prev\_denoised\_ratio\_list} is generated using linear decay:
\begin{lstlisting}
merge_prev_denoised_ratio_list = [
    max_ratio * (1 - t / merge_prev_step)
    for t in range(merge_prev_step)
] + [0] * (num_steps - merge_prev_step)
\end{lstlisting}

\subsection{Issue: Sub-optimal Linear Decay}

Linear decay applies equal blending reduction at each step.
In diffusion models, the early (high-noise) steps contribute coarse structure while late (low-noise) steps refine details.
A linear schedule over-blends during the detail-refinement phase.

\subsection{Proposed Enhancement: Exponential or Cosine Decay}

Use exponential decay so that blending is strong in early steps and decays rapidly:
\begin{equation}
    \lambda(t) = \lambda_{\max} \cdot \exp\left(-\frac{\gamma \, t}{T_{\text{merge}}}\right)
\end{equation}

Or cosine decay:
\begin{equation}
    \lambda(t) = \frac{\lambda_{\max}}{2} \left(1 + \cos\left(\frac{\pi \, t}{T_{\text{merge}}}\right)\right)
\end{equation}

Both schedules reduce interference during the critical detail-refinement phase, improving fine-grained texture quality.

%======================================================================
\section{Logic: Classifier-Free Guidance}
%======================================================================

\subsection{Current Implementation}

CFG is applied uniformly with a constant guidance scale:
\begin{equation}
    \hat{\epsilon} = \epsilon_{\text{uncond}} + s \cdot (\epsilon_{\text{cond}} - \epsilon_{\text{uncond}})
\end{equation}

\subsection{Issue: Over-saturation at Low Noise Levels}

A constant guidance scale $s$ produces good results at high noise levels but tends to cause oversaturated artifacts at low noise levels (later denoising steps) where the signal-to-noise ratio is already high.

\subsection{Proposed Enhancement: Dynamic Guidance Rescaling}

Implement timestep-dependent guidance scaling:
\begin{equation}
    s(t) = s_{\max} \cdot \left(\frac{\bar{\alpha}_t}{1 - \bar{\alpha}_t}\right)^{-\beta}
\end{equation}

where $\beta$ is a small positive constant (e.g.\ $\beta \approx 0.1)$.
This reduces guidance strength as SNR increases.

Alternatively, implement the guidance rescale technique from ``Common Diffusion Noise Schedules and Sample Steps are Flawed'' (Lin et al., 2023):
\begin{equation}
    \hat{\epsilon}_{\text{rescaled}} = \phi \cdot \hat{\epsilon}_{\text{cfg}} + (1-\phi) \cdot \hat{\epsilon}_{\text{cfg}} \cdot \frac{\|\epsilon_{\text{cond}}\|}{\|\hat{\epsilon}_{\text{cfg}}\|}
\end{equation}

This avoids over-exposure without reducing prompt adherence.

%======================================================================
\section{Logic: Tiled VAE Encoding}
%======================================================================

\subsection{Current Implementation}

In \fn{tiled\_vae\_encode\_tensor\_simple()}, the image is split into a fixed $4\!\times\!4$ grid of tiles with a fixed overlap of 32 pixels per side, and the stitched latent is the \emph{average} of overlapping tile latents.

\subsection{Issue: Hard Boundaries Despite Averaging}

Simple averaging at tile boundaries can produce visible seams, because the VAE encoder's receptive field extends beyond the overlap region and encodes context differently at tile edges.

\subsection{Proposed Enhancement: Gaussian-Weighted Blending}

Use Gaussian or cosine-tapered weighting masks for each tile instead of uniform averaging:
\begin{equation}
    W(x, y) = \prod_{d \in \{x, y\}} \begin{cases}
        \frac{1}{2}\left(1 + \cos\left(\frac{\pi (d - d_0)}{O}\right)\right) & \text{if } d \in \text{overlap region} \\
        1 & \text{otherwise}
    \end{cases}
\end{equation}
where $O$ is the overlap width.
This produces smooth transitions:
\begin{equation}
    L_{\text{fused}}(x,y) = \frac{\sum_k W_k(x,y) \cdot L_k(x,y)}{\sum_k W_k(x,y)}
\end{equation}

Additionally, consider making the overlap size and tile count \textbf{adaptive} based on image resolution rather than using fixed constants.

%======================================================================
\section{Logic: Image Preprocessing}
%======================================================================

\subsection{Current Implementation}

In \file{precast\_latent\_utils.py}:
\begin{lstlisting}
transform = transforms.Compose([
    transforms.Resize(image_size),
    transforms.ToTensor(),
    transforms.Normalize(mean=[0.0, 0.0, 0.0], std=[1.0, 1.0, 1.0]),
    transforms.Lambda(lambda x: x * 2.0 - 1.0)
])
\end{lstlisting}

Meanwhile, in \file{tensor\_utils.py}:
\begin{lstlisting}
rgb_img = cv2.resize(rgb_img, (width, height),
                     interpolation=cv2.INTER_LINEAR)
img_tensor = (img_tensor / 255. - 0.5) * 2
\end{lstlisting}

\subsection{Issue 1: Inconsistent Resize Methods}

Two different image loading/resizing pipelines are used (\fn{torchvision.transforms.Resize} vs.\ \fn{cv2.resize}).
These use different interpolation algorithms internally (PIL's Lanczos by default vs.\ OpenCV's bilinear), leading to subtly different pixel values for the same input image.

\subsection{Issue 2: Redundant Normalization}

The \fn{Normalize(mean=[0,...], std=[1,...])} in the transform pipeline is a no-op (identity transform), adding unnecessary computational overhead.

\subsection{Proposed Enhancement}

\begin{enumerate}
    \item \textbf{Unify} all image loading to a single method that uses consistent interpolation (recommend Lanczos/area-based downsampling for higher quality):
\begin{lstlisting}
rgb_img = cv2.resize(rgb_img, (W, H),
                     interpolation=cv2.INTER_AREA)  # for downscale
# or cv2.INTER_LANCZOS4 for upscale
\end{lstlisting}

    \item Remove the redundant \fn{Normalize} transform.
    \item Consider using anti-aliased resizing (e.g., \fn{transforms.Resize(..., antialias=True)}) to prevent aliasing artifacts.
\end{enumerate}

%======================================================================
\section{Logic: Latent Resizing Interpolation}
%======================================================================

\subsection{Current Implementation}

In \file{diffusion\_utils.py}, video latents are resized using \fn{F.interpolate} with either \texttt{nearest} or \texttt{bilinear} mode:
\begin{lstlisting}
upsampled = F.interpolate(video_reshaped,
    size=(target_height, target_width),
    mode=mode, align_corners=align_corners)
\end{lstlisting}

\subsection{Issue: Nearest-Neighbor for Upscaling}

In the sphere panorama pipeline, \fn{resize\_video\_latent(..., mode="nearest")} is used at multiple points for both downscaling and upscaling.
Nearest-neighbor produces blocky latent representations and can introduce high-frequency artifacts in the generated output.

For the upscale path, \texttt{bicubic} is used, which is better, but the upsampled latent is then re-noised without any anti-aliasing.

\subsection{Proposed Enhancement}

\begin{enumerate}
    \item Use \texttt{bilinear} or \texttt{bicubic} for all latent resizing operations, especially upscaling.
    \item Apply a small Gaussian blur to the upsampled latent \emph{before} re-noising to suppress spurious high-frequency content:
\begin{equation}
    L_{\text{smooth}} = G_\sigma * L_{\text{upsampled}}
\end{equation}
where $G_\sigma$ is a Gaussian kernel with $\sigma \approx 0.5$--$1.0$.
\end{enumerate}

%======================================================================
\section{Logic: Shift-Window Theta Offset Cycling}
%======================================================================

\subsection{Current Implementation}

In the sphere panorama sampling:
\begin{lstlisting}
theta_offset = (i % loop_step_theta) * (view_fov // loop_step_theta)
\end{lstlisting}

\subsection{Issue: Integer Division Gaps}

When $\text{view\_fov} = 120$ and $\text{loop\_step\_theta} = 10$, the offset increments by $120 / 10 = 12$ degrees.
But for non-divisible combinations, integer division truncates the step size, leaving \textbf{unvisited angular gaps} in the panoramic coverage.

\subsection{Proposed Enhancement}

Use floating-point division to ensure full coverage:
\begin{equation}
    \theta_{\text{offset}} = (i \bmod N) \cdot \frac{\text{FoV}}{N}
\end{equation}
\begin{lstlisting}
theta_offset = (i % loop_step_theta) * (view_fov / loop_step_theta)
\end{lstlisting}

%======================================================================
\section{Code: Redundant VAE Encoding per Timestep}
%======================================================================

\subsection{Current Implementation}

Inside the denoising loop in \fn{basic\_sample\_shift\_shpere\_panorama()}, the panoramic input image is repeatedly re-encoded by the VAE at every timestep when \fn{paste\_on\_static} is enabled:
\begin{lstlisting}
if paste_on_static and i < total_steps - 1:
    frame_0_latent = self.tiled_vae_encode_image(
        image_path=pano_image_path,
        image_size=(equirect_height, equirect_width))
\end{lstlisting}

\subsection{Issue}

The input image never changes across timesteps, so the VAE encoding produces identical results each time.
This wastes significant GPU computation---the tiled VAE encode over a $1024\!\times\!512$ image is expensive.

\subsection{Proposed Enhancement}

Cache the encoded latent \emph{before} the denoising loop:
\begin{lstlisting}
cached_frame_0_latent = self.tiled_vae_encode_image(
    image_path=pano_image_path,
    image_size=(equirect_height, equirect_width))

for i, t in enumerate(timesteps):
    if paste_on_static and i < total_steps - 1:
        frame_0_latent = cached_frame_0_latent
        ...
\end{lstlisting}

This saves $O(T)$ VAE forward passes, substantially reducing total inference time.

%======================================================================
\section{Code: Redundant \fn{.clone()} Operations}
%======================================================================

\subsection{Current Implementation}

Throughout the codebase, \fn{.clone()} is called extensively, often unnecessarily:
\begin{lstlisting}
# In ring_panorama_tensor_utils.py
view = view_tensor.view(...).clone()

# In scheduler.py
return desired_latent.clone()
\end{lstlisting}

\subsection{Issue}

Many of these clones occur on tensors that are immediately consumed and not referenced again, creating wasteful memory copies.
In a pipeline where latent tensors have shape $(1, 4, 16, 64, 128)$, each clone allocates an additional $\sim$2 MB.

\subsection{Proposed Enhancement}

Audit each \fn{.clone()} and remove those that are not needed for in-place operation safety.
Specifically:
\begin{itemize}
    \item Clones before \fn{scatter\_} or \fn{index\_add\_} operations are necessary---keep them.
    \item Clones on return values that are immediately consumed by the caller---remove them.
    \item In \fn{get\_window\_latent()}, the final \fn{.clone()} can be replaced with returning the concatenated tensor directly.
\end{itemize}

%======================================================================
\section{Code: Duplicate Prompt Encoding}
%======================================================================

\subsection{Current Implementation}

In the sphere panorama pipeline, prompt embeddings are re-computed for \emph{every view} at \emph{every timestep}:
\begin{lstlisting}
text_emb = self.pretrained_t2v.get_learned_conditioning(
                                            [curr_phi_prompt])
img_emb = self.pretrained_t2v.get_image_embeds(
                            batch_imgs=view_condition_image_tensor)
\end{lstlisting}

\subsection{Issue}

The text embedding for each elevation-specific prompt is identical across all timesteps and all theta positions at the same elevation.
These redundant CLIP forward passes waste compute.

\subsection{Proposed Enhancement}

Pre-compute and cache the text embeddings before the denoising loop:
\begin{lstlisting}
phi_text_emb_cache = {}
for phi_angle, phi_prompt in phi_prompt_dict.items():
    phi_text_emb_cache[phi_angle] = \
        self.pretrained_t2v.get_learned_conditioning([phi_prompt])
\end{lstlisting}

Similarly, condition image embeddings can be cached per $(\theta, \phi)$ view since the input panorama does not change.

%======================================================================
\section{Code: Video Saving and Post-processing}
%======================================================================

\subsection{Current Implementation}

In \fn{save\_decoded\_video\_latents()}, video frames are decoded and saved one at a time:
\begin{lstlisting}
for frame_idx in range(decoded_video_latents.shape[2]):
    frame_tensor = decoded_video_latents[:, :, [frame_idx]]
    image = tensor2image(frame_tensor)
    video_frames_img_list.append(image)
\end{lstlisting}

\subsection{Proposed Enhancement}

\begin{enumerate}
    \item Process frames in mini-batches for more efficient GPU utilization.
    \item Use built-in \fn{torchvision.utils.save\_image()} or batched tensor-to-numpy conversion instead of per-frame Python loops.
    \item Add optional temporal smoothing via a 1D Gaussian or bilateral filter along the time dimension before saving, to reduce inter-frame flickering.
\end{enumerate}

%======================================================================
\section{Mathematics: Polar Region Handling}
%======================================================================

\subsection{Current Implementation}

The sphere panorama pipeline treats all elevation angles ($\phi$) uniformly, sampling the same number of theta positions for most elevation angles:
\begin{lstlisting}
phi_theta_dict = {
    90:  [0],
    75:  [360*t//phi_num for t in range(phi_num)],
    ...
    0:   [360*t//phi_num for t in range(phi_num)],
    -75: [360*t//phi_num for t in range(phi_num)],
    -90: [0],
}
\end{lstlisting}

\subsection{Issue: Over-sampling Polar Regions}

At high latitudes (e.g.\ $|\phi| = 75\degree$), the equirectangular projection compresses azimuthal distance.
Sampling 6 theta values at $\phi = 75\degree$ produces views with significant overlap, while at $\phi = 0\degree$ (equator), 6 views may be insufficient.

\subsection{Proposed Enhancement: Latitude-Adaptive Theta Sampling}

Scale the number of theta samples by $\cos(\phi)$ to achieve uniform solid-angle coverage:
\begin{equation}
    N_\theta(\phi) = \max\left(1, \; \left\lfloor N_0 \cdot \cos(\phi) \right\rfloor\right)
\end{equation}
where $N_0$ is the base number of samples at the equator.

This reduces redundant computation at poles and increases coverage at the equator, improving both quality and efficiency.

%======================================================================
\section{Mathematics: Guidance Scale for Multi-View Consistency}
%======================================================================

\subsection{Current Limitation}

Each perspective view is denoised independently with CFG.
There is no mechanism to enforce multi-view consistency beyond the overlap merging.

\subsection{Proposed Enhancement: Multi-View Score Averaging}

Before computing the DDIM step, average the noise predictions from overlapping views:
\begin{equation}
    \hat{\epsilon}_{\text{avg}}(p) = \frac{\sum_{v \in \mathcal{V}(p)} w_v \cdot \hat{\epsilon}_v(p)}{\sum_{v \in \mathcal{V}(p)} w_v}
\end{equation}
where $\mathcal{V}(p)$ is the set of views covering pixel $p$, and $w_v$ accounts for projection distortion (e.g.\ $w_v \propto \cos(\phi_v)$ for sphere-aware weighting).

This is analogous to the ``MultiDiffusion'' approach and would improve global coherence.

%======================================================================
\section{Summary of Proposed Enhancements}
%======================================================================

\begin{table}[H]
\centering
\footnotesize
\renewcommand{\arraystretch}{1.2}
\begin{tabular}{@{}p{0.26\textwidth}p{0.15\textwidth}p{0.10\textwidth}p{0.40\textwidth}@{}}
\toprule
\textbf{Enhancement} & \textbf{Category} & \textbf{Impact} & \textbf{Key Benefit} \\
\midrule
Pixel Grid Centering & Math & Medium & Removes 0.5 px projection bias \\
Longitude Normalization & Math & Medium & Eliminates pixel duplication at 360\textdegree{} seam \\
Nearest-Neighbor Rounding & Math & Low & Reduces average reprojection error \\
Bilinear Splat Vectorization & Code/Math & High & 4--8$\times$ speedup in write-back \\
DDIM-aware Re-noising & Math & High & Better structure preservation in upscaling \\
Smooth Overlap Scheduling & Logic & Medium & Eliminates temporal seam artifacts \\
Cosine/Exp Merge Decay & Logic & Medium & Sharper details in final output \\
Dynamic CFG Guidance & Logic & High & Reduces over-saturation \\
Gaussian Tiled VAE Blending & Logic & Medium & Smoother tile boundaries \\
Unified Image Preprocessing & Code & Low & Consistent pixel values \\
Bicubic Latent Resizing & Logic & Medium & Less blocky upscaling \\
Float Theta Offset & Logic & Low & Eliminates angular gaps \\
Cached VAE Encoding & Code & High & $O(T)$ fewer VAE passes \\
Cached Prompt Embeddings & Code & High & $O(T \!\cdot\! N_\phi \!\cdot\! N_\theta)$ fewer CLIP passes \\
Remove Redundant Clones & Code & Medium & Lower peak memory \\
Latitude-Adaptive Sampling & Math/Logic & High & Uniform coverage, 30\%+ fewer views at poles \\
Multi-View Score Averaging & Math & High & Global panorama coherence \\
\bottomrule
\end{tabular}
\caption{Summary of all proposed enhancements ordered by category.}
\end{table}

%======================================================================
\section{Conclusion}
%======================================================================

The DynamicScaler codebase is a well-structured implementation of a novel panoramic video generation pipeline.
The enhancements proposed in this document span three axes:

\begin{enumerate}
    \item \textbf{Mathematical Precision}: Fixing subtle but impactful geometric projection errors that accumulate across hundreds of view extractions and re-projections.
    \item \textbf{Algorithmic Logic}: Replacing abrupt schedules and uniform strategies with smooth, adaptive alternatives informed by diffusion theory.
    \item \textbf{Code Efficiency}: Caching invariant computations, vectorizing inner loops, and removing redundant memory operations.
\end{enumerate}

Together, these improvements target higher visual quality, fewer seam artifacts, and significantly faster inference---all without requiring architectural changes to the underlying diffusion model.

\end{document}
